\chapter{Concept Testing}
% this is some introduction about the topic
\subsection{Pre-class Activities}
\subsubsection{Activity 1: Define the Purpose of the Concept Test}
\textbf{Objective:}
The objective of this is activity is to better understand the importance of generating questions that you aim to address during the concept testing phase. These questions should provide valuable insights and aid in decision-making
\begin{enumerate}
    \item Draft the questions that your team aim to answer through the concept test

    The typical questions addressed in concept testing are:
    \begin{itemize}
        \item Which of several alternative concepts should we pursue?
        \item How can we improve the concept to better meet customer needs?
        \item Approximately how many units are likely to be sold?
        \item Should we continue with the development of the concept?
    \end{itemize}
    \item Document the question(s)
\end{enumerate}

\subsubsection{Activity 2: Choose a Survey Population} % if available
\textbf{Objective:}
The objective of this activity is to identify the market segment or target audience for your product and select a survey population that closely represents the characteristics of your target market.

\textbf{Procedure:}

\begin{enumerate}
    \item Search for relevant resource to help you understand the specific group of customers (market segment) who are most likely to be interested or benefit from your product.
    \item Once gain an idea, proceed by selecting survey population that closely resembles the characteristics of your target market.
    This population should reflect your potential customers (e.g., age, gender, geographic location, preferences).
    \item Determine the sample size for your survey population.
    \begin{enumerate}
        \item You use free online service such as \href{https://www.qualtrics.com/au/experience-management/research/determine-sample-size}{click me}  to calculate the appropriate sample size for your research
        \item While it is not necessary to survey everyone, it is imperative to take into account this statistical consideration.
    \end{enumerate}
\end{enumerate}

\subsubsection{Activity 3: Choose a Survey Format} % if available
\textbf{Objective:}
The objective of this activity is to select a survey format and develop survey questions that are clear and relevant to the purpose of the concept test.
\begin{enumerate}
    \item Consider the different survey formats available and choose the one that suits your needs. Some common survey formats include:
    \begin{itemize}
        \item Face-to-face interaction
        \item Telephone interviews
        \item Postal mail surveys
        \item Electronic mail surveys
        \item Internet-based surveys
    \end{itemize}
    \item Draft the survey and possible interview questions that will be used during the survey
    \begin{enumerate}
        \item Keep in mind to ensure the questions are clear and relevant to the purpose of the concept test.
    \end{enumerate}
\end{enumerate}

\subsubsection{Activity 4: Communicate the Concept} % if available
\textbf{Objective:}
To objective of this activity to decide on a communication method for presenting the concept to survey participants and ensure that the chosen method allows respondents to understand and evaluate the concept accurately.

\begin{enumerate}
    \item Decide how you will communicate the concept to the survey participants by considering factor such as
    \begin{itemize}
        \item nature of your product
        \item the survey format you have chosen.
    \end{itemize}
    \item Try to come out with one or more means of communication that are suitable for presenting the concept effectively, options including
    \begin{itemize}
        \item Verbal description
        \item Sketches
        \item Photos and renderings
        \item Storyboards
        \item Videos
        \item Simulations
        \item Interactive multimedia
        \item Physical appearance models
    \end{itemize}
    \item Keep in mind that the choice of survey format is closely linked to how you communicate the product concept.
    \begin{enumerate}
        \item Make sure the chosen format allow respondents to understand and evaluate the concept accurately.
        \item Select communication methods that effectively convey the key aspects and attributes of the concept and subsequently facilitating meaningful feedback from participants.
    \end{enumerate}
\end{enumerate}

\subsection{In-class Activities}
\subsubsection{Activity 1: Defining the Purpose of the Concept Test}
\begin{enumerate}
    \item Compile all the findings.
    \begin{enumerate}
        \item Each group member should already individually generate several question(s).
    \end{enumerate}
    \item Discuss with all team members the most advantageous approach:
    \begin{enumerate}
        \item Decide whether to select the best question drafted by one member or combine elements from multiple questions to create a new single question or several questions.
    \end{enumerate}
    \item Once a consensus is reached, draft a question(s) that reflects the collective agreement of the team.
\end{enumerate}

\subsubsection{Activity 2: Choosing a Survey Population} % if available
\begin{enumerate}
    \item Compile all the findings.
    \begin{enumerate}
        \item Each group member should already individually generate several market segment and survey population.
    \end{enumerate}
    \item Discuss with all team members whether it is more advantageous to
    \begin{enumerate}
        \item select the best market segment and survey population drafted by one member or
        \item combine elements from market segments and survey populations drafted by all team members.
    \end{enumerate}
    \item Once a consensus is reached, draft a market segment and survey population that reflects the collective agreement of the team.
\end{enumerate}

\subsubsection{Activity 3: Choosing a Survey Format} % if available
\begin{enumerate}
    \item Compile all the findings.
    \begin{enumerate}
        \item Each group member should already individually draft their justification about the best survey format for the chosen product.
    \end{enumerate}
    \item Discuss the advantage and disadvantage the proposal given by each member.
    \item Once a consensus is reached, pick and documented the suitable survey format.
\end{enumerate}

\subsubsection{Activity 4: Communicating the Concept} % if available
\begin{enumerate}
    \item Discuss with all team members whether it is more advantageous to select the best communication method proposed by one member or combine elements from multiple communication methods.
    \item Once a consensus is reached, propose a communication method that reflects the collective agreement of the team.
\end{enumerate}